% =========================================================
%  frontmatter/copyright.tex
%  Halaman hak cipta dan keterangan penggunaan
% =========================================================

\chapter*{Hak Cipta dan Keterangan Penggunaan}
\addcontentsline{toc}{chapter}{Hak Cipta dan Keterangan Penggunaan}
\thispagestyle{plain}

\vspace{0.5cm}

{\bfseries \BookTitle\par}
{\itshape \BookSubtitle\par}
\vspace{0.4cm}

Hak cipta \copyright\ \BookYear\ pada \NamaDosen, Program Studi \NamaProdi,
\NamaFakultas, \NamaPT.

\vspace{0.6cm}

{\bfseries Penerbit\par}
\NamaPT\par
\NamaFakultas\par
Program Studi \NamaProdi

\vspace{0.9cm}

{\bfseries Keterangan Penggunaan\par}

Buku ajar ini disusun untuk keperluan pembelajaran mata kuliah \NamaMK\
(\KodeMK) di lingkungan \NamaPT. Penggandaan, penyalinan, atau penyebaran
sebagian maupun seluruh isi buku untuk keperluan pendidikan nonkomersial
diperbolehkan dengan menyebutkan sumbernya. Penggunaan untuk keperluan
komersial memerlukan izin tertulis dari penulis.

\vspace{0.6cm}

{\bfseries Keterangan atas Materi Pihak Ketiga\par}

Buku ini memuat rujukan dan, pada beberapa bagian, gambar pendukung yang
berasal dari sumber pihak ketiga. Setiap gambar yang bukan karya penulis
diberi keterangan sumber di bawah gambar yang bersangkutan dengan nomor
rujukan \mbox{[R1]}--\mbox{[R35]} sesuai daftar pustaka. Seluruh merek
dagang, nama produk, dan logo yang disebutkan dalam buku ini merupakan
milik pemegang haknya masing-masing, termasuk Python, Pygame, dan Unity
yang digunakan semata-mata sebagai perkakas pembelajaran.

\vspace{0.6cm}

{\bfseries Penafian\par}

Buku ajar ini mengikuti Rencana Pembelajaran Semester (RPS) mata kuliah
\NamaMK\ yang berlaku pada Program Studi \NamaProdi\ \NamaPT. Perubahan
RPS, versi perkakas, atau tautan sumber belajar di luar kendali penulis
dapat menyebabkan sebagian isi buku perlu disesuaikan pada edisi
berikutnya. Tautan daring pada daftar pustaka ditelusuri pada September
2026 dan perlu diverifikasi ulang ketika buku digunakan.

\vspace{1.0cm}

\begin{tabular}{@{}L{2.6cm}L{9.6cm}@{}}
\textbf{Edisi} & \BookEdition \\
\textbf{Tahun} & \BookYear \\
\textbf{Penulis} & \BookAuthor \\
\textbf{Mata kuliah} & \NamaMK\ (\KodeMK), \BobotSKS \\
\end{tabular}

\clearpage
